\documentclass[11pt, a4paper]{article}

\usepackage[T1]{fontenc}
\usepackage[utf8]{inputenc}
\usepackage{tgpagella}
\usepackage[spanish, es-noshorthands]{babel}
\usepackage{amsmath, amssymb}
\usepackage[most]{tcolorbox}
\usepackage[american]{circuitikz}
\usepackage[margin=2.5cm]{geometry}
\usepackage{fancyhdr}

\pagestyle{fancy}
\fancyhf{}
\setlength{\headheight}{16pt}
\lhead{\textbf{UCB Sede Tarija} | Ing. Mecatrónica}
\rhead{IMT-221 | Diagnóstico AC S06-3}
\renewcommand{\headrulewidth}{0.4pt}
\definecolor{ucbblue}{RGB}{0, 51, 102}

\begin{document}

\begin{tcolorbox}[colback=ucbblue!5, colframe=ucbblue, title=\textbf{Consolidación de Pequeña Señal}]
    Tiempo estimado: 7-10 minutos. Evaluación de autonomía en parámetros intrínsecos.
\end{tcolorbox}

\section*{Ejercicio 1: Emisor Común Aislado}
Considera el circuito mostrado, con punto de operación DC ya determinado. \\
\textbf{Datos asignados:} $I_{CQ}=1.0\,\text{mA}$, $\beta=100$, $R_C=4.7\,\text{k}\Omega$, $T \approx 300\,\text{K}$ ($V_T=25.85\,\text{mV}$). 

\begin{center}
\begin{circuitikz}[scale=1, transform shape, thick]
    \draw (0,0) node[npn](Q){};
    \draw (Q.E) node[ground]{};
    \draw (Q.C) to[R, l_={$R_C$}] (0,2) -- (0,2.5) node[above]{$+V_{CC}$};
    \draw (Q.B) -- (-2,0) node[left]{$v_i$};
    \draw (Q.C) to[short, -o] (1,0.77) node[right]{$v_o$};
\end{circuitikz}
\end{center}

\textbf{Calcula y responde secuencialmente:}
\begin{enumerate}
    \item Calcula la transconductancia $g_m$. \vspace{1.5cm}
    \item Calcula la resistencia de base en pequeña señal $r_\pi$. \vspace{1.5cm}
    \item Dibuja el modelo Híbrido-$\pi$ equivalente del transistor, conectando los terminales B, C y E, y ubicando correctamente la Tierra AC. \vspace{3cm}
    \item Estima la ganancia de etapa sin carga $A_{v0} = -g_m R_C$. \vspace{1cm}
    \item Nombra al menos una suposición técnica que sustente el cálculo anterior.
\end{enumerate}

\end{document}
